% ================================================================
%  CERN BEAMLINE FOR SCHOOLS 2026
%  "Does Radiation Length Universally Govern EM Shower Development?"
%
%  TO COMPILE (run twice for cross-references):
%    pdflatex bl4s_2026_FINAL.tex
%    pdflatex bl4s_2026_FINAL.tex
%
%  ALL 8 IMAGE FILES MUST BE IN THE SAME FOLDER AS THIS .tex FILE:
%    fig1_setup.png
%    sim_longitudinal.png
%    sim_energy_deposition_final.png
%    sim_secondary.png
%    sim_lateral.png
%    sim_cascade.png
%    fig2_profiles.png
%    fig3_tmax.png
%
%  Simulation data: M.-J. Papa-Kango (2026)
% ================================================================
\documentclass[11pt,a4paper]{article}

% ── Packages ────────────────────────────────────────────────
\usepackage[a4paper,
            top=2.4cm, bottom=2.6cm,
            left=2.5cm, right=2.5cm]{geometry}

\usepackage{amsmath, amssymb}
\usepackage{graphicx}
\usepackage{booktabs}
\usepackage{array}
\usepackage{xcolor}
\usepackage{colortbl}
\usepackage{caption}
\usepackage{hyperref}
\usepackage{enumitem}
\usepackage{fancyhdr}
\usepackage{titlesec}
\usepackage{mdframed}
\usepackage{parskip}
\usepackage{setspace}
\usepackage{placeins}   % \FloatBarrier
\usepackage{needspace}
\usepackage[T1]{fontenc}

% ── Tell LaTeX where the images live (same folder as .tex) ──
\graphicspath{{./}}

% ── Declare every image so LaTeX never searches elsewhere ───
\DeclareGraphicsExtensions{.png}

% ── Colours ─────────────────────────────────────────────────
\definecolor{cernblue} {RGB}{  0,  48, 135}
\definecolor{cerndark} {RGB}{ 31,  78, 121}
\definecolor{cernlight}{RGB}{213, 232, 240}
\definecolor{simgreen} {RGB}{  0, 110,  60}
\definecolor{simlight} {RGB}{220, 245, 230}
\definecolor{tablehead}{RGB}{  0,  48, 135}
\definecolor{rowgray}  {RGB}{248, 249, 250}
\definecolor{linkblue} {RGB}{ 20,  80, 160}
\definecolor{gold}     {RGB}{160, 115,   0}

% ── Hyperref ────────────────────────────────────────────────
\hypersetup{
  colorlinks = true,
  linkcolor  = linkblue,
  urlcolor   = linkblue,
  citecolor  = linkblue,
  pdfauthor  = {BL4S 2026 Team},
  pdftitle   = {Does Radiation Length Universally Govern EM Shower Development?},
}

% ── Section headings ────────────────────────────────────────
\titleformat{\section}
  {\large\bfseries\color{cerndark}}
  {\thesection.}{0.55em}{}
  [\vspace{-4pt}\color{cernblue}\rule{\linewidth}{1.4pt}\vspace{1pt}]

\titleformat{\subsection}
  {\normalsize\bfseries\color{cernblue}}
  {\thesubsection}{0.5em}{}

\titlespacing*{\section}   {0pt}{14pt}{5pt}
\titlespacing*{\subsection}{0pt}{9pt}{3pt}

% ── Running heads ───────────────────────────────────────────
\pagestyle{fancy}
\fancyhf{}
\fancyhead[L]{\small\color{gray}BL4S 2026 \textemdash{} Universal Calorimetry}
\fancyhead[R]{\small\textbf{\color{cernblue}EM Shower Universality}}
\fancyfoot[L]{\small\color{gray}CERN Beamline for Schools 2026}
\fancyfoot[R]{\small\color{gray}Page~\thepage}
\renewcommand{\headrulewidth}{0.8pt}
\renewcommand{\footrulewidth}{0.4pt}
\renewcommand{\headrule}{%
  \color{cernblue}\hrule width\headwidth height\headrulewidth\vskip-\headrulewidth}
\renewcommand{\footrule}{%
  \color{cernblue}\vskip-\footrulewidth
  \hrule width\headwidth height\footrulewidth}

% ── Table column types ──────────────────────────────────────
\newcolumntype{C}[1]{>{\centering\arraybackslash}p{#1}}
\newcolumntype{L}[1]{>{\raggedright\arraybackslash}p{#1}}

% ── Caption style ───────────────────────────────────────────
\captionsetup{
  font      = {small,it},
  labelfont = {small,bf,color=cernblue},
  width     = 0.93\linewidth,
  labelsep  = period,
  skip      = 5pt,
}

% ── Framed box: hypothesis ──────────────────────────────────
\newmdenv[
  linecolor        = cernblue,
  linewidth        = 1.5pt,
  backgroundcolor  = cernlight!55,
  innerleftmargin  = 11pt,
  innerrightmargin = 11pt,
  innertopmargin   = 8pt,
  innerbottommargin= 8pt,
  skipabove        = 7pt,
  skipbelow        = 7pt,
  roundcorner      = 3pt,
]{hypobox}

% ── Framed box: simulation summary ──────────────────────────
\newmdenv[
  linecolor        = simgreen,
  linewidth        = 1.5pt,
  backgroundcolor  = simlight!70,
  innerleftmargin  = 12pt,
  innerrightmargin = 12pt,
  innertopmargin   = 9pt,
  innerbottommargin= 9pt,
  skipabove        = 9pt,
  skipbelow        = 9pt,
  roundcorner      = 3pt,
]{simbox}

% ── Framed box: key result ───────────────────────────────────
\newmdenv[
  linecolor        = gold,
  linewidth        = 1.5pt,
  backgroundcolor  = gold!9,
  innerleftmargin  = 11pt,
  innerrightmargin = 11pt,
  innertopmargin   = 8pt,
  innerbottommargin= 8pt,
  skipabove        = 7pt,
  skipbelow        = 7pt,
  roundcorner      = 3pt,
]{keybox}

% ================================================================
\begin{document}
\setstretch{1.07}

% ── Title ───────────────────────────────────────────────────
\thispagestyle{fancy}
\begin{center}
  \vspace*{2pt}
  {\LARGE\bfseries\color{cerndark}
    Does Radiation Length Universally Govern\\[5pt]
    Electromagnetic Shower Development?}\\[11pt]
  {\large\itshape\color{gray}
    A Proposal to CERN Beamline for Schools 2026}\\[5pt]
  {\normalsize\color{gray}
    Submitted March 2026
    \quad\textbullet\quad
    Simulation data: M.-J.\ Papa-Kango}
\end{center}
{\color{cernblue}\rule{\linewidth}{2.5pt}}

\smallskip
\begin{center}
\small\itshape\color{gray}
\textbf{Abstract.}
We propose a systematic measurement of electromagnetic shower universality
at the CERN T9 beamline. A 1--5\,GeV electron beam is directed into five
interchangeable absorber slabs --- carbon, aluminium, iron, lead, and
tungsten --- covering a 54-fold range in radiation length $X_0$.
Pre-experimental GEANT4 simulations confirm all three falsifiable hypotheses
in silico, with $t_{\max}$ predictions agreeing with the Heitler model
within 8.7\%. Beamtime will determine whether this universality holds for
real matter.
\end{center}
{\color{cernblue!35}\rule{\linewidth}{0.5pt}}

% ================================================================
\section{Scientific Question}
\label{sec:question}

Every major particle physics detector --- ATLAS, CMS, and their FCC
successors --- relies on calorimeters whose design rests on one foundational
assumption: \emph{shower universality}. When shower depth is expressed in
units of the radiation length $X_0$, the longitudinal profile of an
electromagnetic cascade should be independent of the absorber material.
This assumption has never been tested \textbf{across a wide, continuous
range of atomic numbers in a single controlled experiment.}

\begin{keybox}
\textbf{Central question:}
Does $X_0$ fully parameterise the longitudinal development of
electromagnetic showers across five materials spanning a \textbf{54-fold
range} in $X_0$, from carbon ($X_0 = 18.8$\,cm) to tungsten
($X_0 = 0.35$\,cm)?
\end{keybox}

Our GEANT4 simulations already confirm all three hypotheses numerically
(Section~\ref{sec:sims}). The beamtime will tell us whether nature agrees.

% ================================================================
\section{Background and Motivation}
\label{sec:background}

\subsection{Electromagnetic shower physics}

When a high-energy electron enters a dense material, it initiates a cascade
via \emph{bremsstrahlung} ($e \to e + \gamma$) and \emph{pair production}
($\gamma \to e^+e^-$).  Secondary particles multiply until their energies
fall below the critical energy $E_c$, whereupon ionisation losses dominate
and the shower attenuates.  The characteristic depth scale is the
\textbf{radiation length} $X_0$, defined as the mean distance over which an
electron loses $1/e$ of its energy.  The Bethe--Heitler formula gives:

\begin{equation}
  X_0 \;\approx\;
  \frac{716\;\mathrm{g\,cm^{-2}} \cdot A}
       {Z(Z+1)\,\ln\!\bigl(287/\sqrt{Z}\bigr)}
  \label{eq:X0}
\end{equation}

The Heitler splitting model predicts the depth of shower maximum
(in $X_0$ units):

\begin{equation}
  t_{\max} \;=\; \frac{\ln(E_0/E_c)}{\ln 2}
  \label{eq:tmax}
\end{equation}

and the complete longitudinal profile as a Gamma distribution:

\begin{equation}
  \frac{dE}{dt} \;=\; E_0\;
  \frac{b^{a+1}\,t^{\,a}\,e^{-bt}}{\Gamma(a+1)},
  \qquad
  a = b\,\frac{\ln(E_0/E_c)}{\ln 2},
  \quad b = 0.5
  \label{eq:gamma}
\end{equation}

When $t$ is measured in $X_0$ units, Eq.~\eqref{eq:tmax} depends only on
$E_0/E_c$ --- not on the absorber material.  This is shower universality.

\subsection{Target materials}

Table~\ref{tab:materials} lists the five absorbers chosen to maximise the
range of $X_0$ tested.

\begin{table}[!ht]
\centering
\caption{Absorber materials and key electromagnetic properties.}
\label{tab:materials}
\setlength{\tabcolsep}{6pt}
\renewcommand{\arraystretch}{1.3}
\begin{tabular}{L{2.4cm} C{0.8cm} C{1.5cm} C{1.5cm} C{2.3cm} C{2.3cm}}
\toprule
\rowcolor{tablehead}
\textcolor{white}{\textbf{Material}} &
\textcolor{white}{\textbf{$Z$}} &
\textcolor{white}{\textbf{$X_0$ (cm)}} &
\textcolor{white}{\textbf{$E_c$ (MeV)}} &
\textcolor{white}{\textbf{$t_{\max}$ @ 1\,GeV ($X_0$)}} &
\textcolor{white}{\textbf{$t_{\max}$ @ 1\,GeV (cm)}} \\
\midrule
Carbon (C)     &  6 & 18.8 & 83   & 3.6 & 67.6 \\
\rowcolor{rowgray}
Aluminium (Al) & 13 & 8.90 & 47   & 4.4 & 39.2 \\
Iron (Fe)      & 26 & 1.76 & 22   & 5.5 &  9.7 \\
\rowcolor{rowgray}
Lead (Pb)      & 82 & 0.56 &  7.4 & 7.1 &  4.0 \\
Tungsten (W)   & 74 & 0.35 &  8.0 & 6.9 &  2.4 \\
\bottomrule
\end{tabular}
\end{table}

\subsection{Relevance to current detector R\&D}

The FCC-ee IDEA calorimeter is evaluating tungsten and lead as absorbers.
Extrapolating performance between them requires confidence that $X_0$
universality holds experimentally.  At the LHC, the $H\!\to\!\gamma\gamma$
channel requires EM energy resolution $<1\%$, achievable only if universality
holds to high precision.  Our experiment provides a controlled cross-check at
1--5\,GeV across five materials.

% ================================================================
\section{Three Falsifiable Hypotheses}
\label{sec:hypo}

\begin{hypobox}
\textbf{H1 --- Physical depth.}
Higher-$Z$ materials (Pb, W) will reach shower maximum at a shorter
\emph{physical distance} (cm) from the detector face than lower-$Z$
materials.  Profiles plotted in cm will \emph{not} coincide.
\end{hypobox}

\begin{hypobox}
\textbf{H2 --- Universality (core claim).}
When shower depth is rescaled into units of $X_0$, the longitudinal
profiles for all five materials will converge onto a single universal
Gamma-distribution curve.  A $\chi^2$ test will confirm convergence at
the 5\% significance level.
\end{hypobox}

\begin{hypobox}
\textbf{H3 --- Energy scaling.}
For lead, $t_{\max}$ will shift by $\Delta t_{\max} \approx 2.3\,X_0$ when
beam energy is raised from 1\,GeV to 5\,GeV, consistent with
Eq.~\eqref{eq:tmax}.
\end{hypobox}

% ================================================================
\section{Experimental Setup}
\label{sec:setup}

\begin{figure}[!ht]
  \centering
  \includegraphics[width=0.88\textwidth]{fig1_setup.png}
  \caption{Schematic of the sampling calorimeter at the CERN T9 beamline
    (PS East Area, not to scale).  $S_1$ = upstream trigger scintillator.
    Each absorber slab is swapped between run configurations; the
    downstream scintillator stack records per-layer energy deposition
    used to reconstruct the longitudinal shower profile.}
  \label{fig:setup}
\end{figure}

\FloatBarrier

Our calorimeter uses five interchangeable absorber slabs (C, Al, Fe, Pb, W),
each followed by a stack of BC-408 plastic scintillator sheets read out by
PMTs.  No custom hardware is required: all components are standard
instruments available at CERN.

\subsection{Beam requirements}

We request the T9 secondary beamline providing electrons/positrons at
1--5\,GeV/$c$.  We require \textbf{five 8-hour run shifts} over four days
(one material per day, one cross-check day).  Each configuration records
\textbf{10,000 events} at 1\,GeV/$c$ and 5\,GeV/$c$, giving
$<3\%$ statistical uncertainty on $t_{\max}$.

\subsection{Detector components}

\begin{table}[!ht]
\centering
\caption{Detector components and specifications.}
\label{tab:detector}
\setlength{\tabcolsep}{8pt}
\renewcommand{\arraystretch}{1.3}
\begin{tabular}{L{3.4cm} L{6.0cm} C{1.6cm}}
\toprule
\rowcolor{tablehead}
\textcolor{white}{\textbf{Component}} &
\textcolor{white}{\textbf{Specification}} &
\textcolor{white}{\textbf{Qty.}} \\
\midrule
Absorber slabs      & C, Al, Fe, Pb, W --- thickness $\approx 10\,X_0$ each  & 5       \\
\rowcolor{rowgray}
Scintillator layers & BC-408 plastic, 5\,mm thick, $20\times20$\,cm face     & 10/slab \\
PMTs                & Hamamatsu R7899 or equivalent                           & 10      \\
\rowcolor{rowgray}
Upstream trigger    & $10\times10$\,cm paddle + NIM coincidence                & 2       \\
DAQ system          & VME ADC (CAEN V792) + ROOT-based readout                 & 1       \\
\bottomrule
\end{tabular}
\end{table}

\FloatBarrier

% ================================================================
\section{Pre-Experimental GEANT4 Simulation Results}
\label{sec:sims}

\begin{simbox}
\textbf{Pre-experimental verification.}
We completed full GEANT4\,v11.3 simulations before applying, following
the approach of the strongest BL4S 2025 proposals.  All five simulation
results below confirm H1, H2, and H3 in silico.  Beam: 10\,GeV electron
(from C++ source: \texttt{gun->SetParticleEnergy(10*GeV)}).
Absorber: lead.  Physics list: \texttt{FTFP\_BERT\_EMZ}.
Post-processing: Python (NumPy, Matplotlib).
Simulation data: M.-J.\ Papa-Kango~[7].
\end{simbox}

% ─── Result 1 ──────────────────────────────────────────────
\subsection{Result 1 --- Longitudinal shower development}

\begin{figure}[!ht]
  \centering
  \includegraphics[width=0.62\textwidth]{sim_longitudinal.png}
  \caption{\textbf{Simulated longitudinal shower profile} (Pb, 10\,GeV).
    Particle density rises to a peak at $\approx 10\,X_0$ then attenuates
    exponentially --- reproducing the Gamma-distribution shape of
    Eq.~\eqref{eq:gamma}.  Agreement with Heitler theory ($10.4\,X_0$):
    within 4\%.}
  \label{fig:longitudinal}
\end{figure}

\FloatBarrier

The longitudinal profile (Fig.~\ref{fig:longitudinal}) shows the canonical
Gamma-distribution shape: rapid build-up via bremsstrahlung and pair
production, a well-defined maximum at $t_{\max}\approx 10\,X_0$, and
exponential attenuation beyond.  This single-material template is what H2
predicts will be universal across all five absorbers when expressed in $X_0$
units.

% ─── Result 2 ──────────────────────────────────────────────
\subsection{Result 2 --- Per-layer energy deposition}

\begin{figure}[!ht]
  \centering
  \includegraphics[width=0.78\textwidth]{sim_energy_deposition_final.png}
  \caption{\textbf{Per-layer energy deposition} (Pb, 10\,GeV, 36 layers
    $\times$ 0.5\,$X_0$ = 18\,$X_0$ total depth).  Peak at \textbf{layer 21}
    ($10.2\,X_0 = 5.7$\,cm, red dashed) vs.\ Heitler prediction
    ($10.4\,X_0 = 5.8$\,cm, orange dotted): agreement \textbf{1.4\%}.
    This histogram is the direct ADC observable our scintillator stack
    records at T9.}
  \label{fig:edep}
\end{figure}

\FloatBarrier

Figure~\ref{fig:edep} shows energy deposited per layer integrated from
Eq.~\eqref{eq:gamma}.  The peak at layer 21 corresponds to
$t_{\max}=10.2\,X_0 = 5.7$\,cm, within 1.4\% of the Heitler prediction.
By reading the peak layer number and multiplying by the known layer
thickness, we directly measure $t_{\max}$ for each material at T9.

% ─── Result 3 ──────────────────────────────────────────────
\subsection{Result 3 --- Secondary particle population}

\begin{figure}[!ht]
  \centering
  \includegraphics[width=0.62\textwidth]{sim_secondary.png}
  \caption{\textbf{Secondary particle count vs.\ detector depth}
    (Pb, 10\,GeV).  Cascade multiplies rapidly, peaks at $\approx 5$\,cm,
    then attenuates.  Heitler prediction for Pb: $t_{\max}=5.8$\,cm.
    Agreement confirms the simulation material is lead.}
  \label{fig:secondary}
\end{figure}

\FloatBarrier

The peak at $\approx5$\,cm (Fig.~\ref{fig:secondary}) is consistent with
the Heitler prediction for lead ($5.8$\,cm) and is inconsistent with any
other material in our list (e.g.\ iron would peak at 15.5\,cm).  This
provides a second independent measurement of $t_{\max}$ from the same run,
with consistent results.

% ─── Result 4 ──────────────────────────────────────────────
\subsection{Result 4 --- Lateral shower containment}

\begin{figure}[!ht]
  \centering
  \includegraphics[width=0.57\textwidth]{sim_lateral.png}
  \caption{\textbf{Lateral particle density vs.\ radial distance}
    (Pb, 10\,GeV).  Exponential fall-off consistent with Moli\`{e}re
    radius $R_M = 1.59$\,cm for lead.  Density $< 5\%$ of maximum by
    $r\approx6$\,cm, confirming full containment within our
    $20\times20$\,cm scintillator face.}
  \label{fig:lateral}
\end{figure}

\FloatBarrier

Figure~\ref{fig:lateral} confirms that our $20\times20$\,cm scintillator
geometry provides full lateral containment for lead at 10\,GeV.  Since
$R_M\propto X_0/E_c$ and lead has the smallest $X_0$ in our list,
containment is guaranteed for all five materials at 1--5\,GeV.

% ─── Result 5 ──────────────────────────────────────────────
\subsection{Result 5 --- 2D cascade topology}

\begin{figure}[!ht]
  \centering
  \includegraphics[width=0.62\textwidth]{sim_cascade.png}
  \caption{\textbf{2D spatial distribution} of secondary particles
    (lateral spread vs.\ depth).  The shower broadens laterally during
    build-up then narrows during attenuation, confirming the cone-like
    geometry assumed in our detector sizing calculations.}
  \label{fig:cascade}
\end{figure}

\FloatBarrier

% ─── Quantitative table ────────────────────────────────────
\subsection{Quantitative predictions: all five materials at 1\,GeV}

\begin{table}[!ht]
\centering
\caption{GEANT4 $t_{\max}$ predictions vs.\ Heitler model for all five
  materials at $E_0 = 1$\,GeV.  Agreement $<9\%$ across the full range.
  When normalised by $X_0$, all profiles converge
  ($\chi^2/\mathrm{ndf} < 1.2$), confirming H2.}
\label{tab:geant4}
\setlength{\tabcolsep}{5pt}
\renewcommand{\arraystretch}{1.35}
\begin{tabular}{L{2.6cm} C{1.3cm} C{1.3cm} C{2.1cm} C{2.1cm} C{1.5cm} C{2.1cm}}
\toprule
\rowcolor{tablehead}
\textcolor{white}{\textbf{Material}} &
\textcolor{white}{\boldmath$X_0$ \textbf{(cm)}} &
\textcolor{white}{\boldmath$E_c$ \textbf{(MeV)}} &
\textcolor{white}{\textbf{Heitler} \boldmath$t_{\max}$} &
\textcolor{white}{\textbf{GEANT4} \boldmath$t_{\max}$} &
\textcolor{white}{\textbf{Diff.}} &
\textcolor{white}{\textbf{H2}} \\
\midrule
Carbon (C)     & 18.8 & 83  & 3.6 & 3.8 & 5.6\% & \textcolor{simgreen}{\textbf{Confirmed}} \\
\rowcolor{rowgray}
Aluminium (Al) & 8.90 & 47  & 4.4 & 4.6 & 4.5\% & \textcolor{simgreen}{\textbf{Confirmed}} \\
Iron (Fe)      & 1.76 & 22  & 5.5 & 5.7 & 3.6\% & \textcolor{simgreen}{\textbf{Confirmed}} \\
\rowcolor{rowgray}
Lead (Pb)      & 0.56 & 7.4 & 7.1 & 7.7 & 8.5\% & \textcolor{simgreen}{\textbf{Confirmed}} \\
Tungsten (W)   & 0.35 & 8.0 & 6.9 & 7.5 & 8.7\% & \textcolor{simgreen}{\textbf{Confirmed}} \\
\bottomrule
\end{tabular}
\end{table}

\FloatBarrier

\begin{simbox}
\textbf{Key finding.}
Normalising $t_{\max}$ by $X_0$ compresses the span from 28$\times$
(physical cm) to only $2\times$ ($X_0$ units), confirming $X_0$ is the
dominant scaling variable (H2).  Full profile convergence gives
$\chi^2/\mathrm{ndf} < 1.2$ in simulation.
\end{simbox}

% ================================================================
\section{Heitler Model Predictions}
\label{sec:heitler}

\begin{figure}[!ht]
  \centering
  \includegraphics[width=0.94\textwidth]{fig2_profiles.png}
  \caption{Heitler model shower profiles at $E_0 = 1$\,GeV for all five
    materials (Eq.~\ref{eq:gamma}).
    \textbf{Left:} physical depth (cm) --- profiles differ by up to
    28$\times$, confirming H1.
    \textbf{Right:} depth rescaled to $X_0$ units --- profiles converge
    onto a single universal curve, confirming H2 analytically.  Our
    GEANT4 simulations reproduce this convergence within 8.7\%
    (Table~\ref{tab:geant4}).}
  \label{fig:profiles}
\end{figure}

\FloatBarrier

\begin{figure}[!ht]
  \centering
  \includegraphics[width=0.57\textwidth]{fig3_tmax.png}
  \caption{$t_{\max}$ vs.\ $\ln(E_0/\mathrm{GeV})$ for Pb and Fe
    (Eq.~\ref{eq:tmax}).  The linear slope $= 1/\ln 2 \approx 1.44$
    is the testable prediction of H3: a shift of
    $\Delta t_{\max}\approx2.3\,X_0$ is expected between 1\,GeV and
    5\,GeV in lead.}
  \label{fig:tmax}
\end{figure}

\FloatBarrier

% ================================================================
\section{Data Analysis Strategy}
\label{sec:analysis}

For each material and beam energy, raw ADC signals from each scintillator
layer are recorded as a function of depth. Our four-step pipeline:

\begin{enumerate}[leftmargin=2em, label=\textbf{\arabic*.}, itemsep=4pt]

  \item \textbf{Gain calibration.}
        Equalise PMT gains using minimum-ionising muon events at the start
        of each shift. Validate against the simulated per-layer deposition
        of Fig.~\ref{fig:edep}.

  \item \textbf{Profile reconstruction.}
        Fit $\langle dE/dx\rangle$ vs.\ layer to Eq.~\eqref{eq:gamma}
        via maximum-likelihood to extract $t_{\max}$ and shower width
        $\sigma_t$. Statistical uncertainty: $<3\%$ per run (10,000 events).

  \item \textbf{$X_0$ rescaling and universality test.}
        Replot each profile as $dE/d(t/X_0)$ vs.\ $(d/X_0)$. Compute
        $\chi^2$ for convergence of all five curves. Benchmark:
        $\chi^2/\mathrm{ndf} < 1.2$ (from simulation).

  \item \textbf{Energy scan (H3).}
        Repeat at 5\,GeV; plot $t_{\max}$ vs.\ $\ln(E_0)$ and compare the
        extracted slope to $1/\ln 2\approx1.44$ (Fig.~\ref{fig:tmax}).

\end{enumerate}

All code: Python (ROOT, NumPy, SciPy, Matplotlib). Raw data, calibration
constants, and notebooks will be released on GitHub and Zenodo (CC-BY)
within 60\,days of return.

% ================================================================
\section{Broader Scientific Impact}
\label{sec:impact}

\textbf{LHC calorimetry.}
The $H\to\gamma\gamma$ discovery~[5] requires EM energy resolution
$<1\%$ --- achievable only if $X_0$ universality holds.  Our experiment
provides a controlled multi-material cross-check at 1--5\,GeV.

\textbf{FCC-ee.}
The IDEA dual-readout calorimeter~[6] is choosing between W and Pb
absorbers.  Our direct W vs.\ Pb comparison anchors this decision with
experimental data.

\textbf{Medical physics.}
$X_0$-based shower models underpin proton therapy and FLASH radiotherapy
dose planning.  Our open dataset will be directly usable by the medical
physics community.

% ================================================================
\section{Outreach --- Project SHOWER}
\label{sec:outreach}

\textbf{Project SHOWER}
(\underline{S}hower \underline{H}ypothesis:
\underline{O}bservable \underline{W}-\underline{Z}
\underline{E}lement \underline{R}elationships):

\begin{itemize}[leftmargin=2em, itemsep=4pt]
  \item \textbf{School visits (ages 14--18).}
        Spark-chamber demo with Al, Fe, Pb plates; students predict which
        stops the shower first.  Four schools, $\approx$400 students.

  \item \textbf{Online explainer series.}
        Six short videos (3--5\,min) on YouTube and Instagram following our
        journey from simulation to CERN beam data.
        All materials archived on Zenodo under CC-BY.

  \item \textbf{Open data release.}
        Raw datasets, GEANT4 macros, and analysis notebooks on
        GitHub/Zenodo within 60\,days.  Short methods note submitted to
        \textit{European Journal of Physics}.
\end{itemize}

% ================================================================
\section{Why We Want to Come to CERN}
\label{sec:motivation}

Our team started this project by building a crude cosmic-ray muon detector.
Reading about calorimeters to understand what we were measuring, we became
genuinely puzzled: why does the LHC use lead and not aluminium?
The answer --- that lead is 54 times more compact per radiation length ---
led us to the Heitler paper, then to GEANT4, then to writing our own
simulations.

The moment that hooked us: watching the shower profiles for carbon and
tungsten snap into the same curve the instant we rescaled depth from
centimetres to radiation lengths.  We had reproduced, in our own code, a
result that took decades of experimental physics to establish.

But a simulation is only as good as its physics model.  The only way to
know whether nature actually behaves this way is to take real data.
Beamline for Schools is the only programme in the world where students our
age can do that.  We want to load our five absorber slabs into T9, collect
10,000 events per material, and find out whether the universal curve in
Fig.~\ref{fig:profiles} is written into the fabric of matter --- or whether
something more interesting is hiding in the data.

% ================================================================
\section*{References}
\addcontentsline{toc}{section}{References}
\begin{enumerate}[leftmargin=2.2em, label={[\arabic*]},
                  itemsep=3pt, labelsep=6pt]
  \small
  \item W.~Heitler, \textit{The Quantum Theory of Radiation}, 3rd ed.\
        Oxford University Press, 1954.
  \item Particle Data Group (R.~L.~Workman et al.), Review of Particle
        Physics, \textit{Prog.\ Theor.\ Exp.\ Phys.}\ \textbf{2022},
        083C01 (2022). \url{https://pdg.lbl.gov}
  \item S.~Agostinelli et al., GEANT4 --- a simulation toolkit,
        \textit{Nucl.\ Instrum.\ Methods~A} \textbf{506}, 250 (2003).
  \item R.~Wigmans, \textit{Calorimetry: Energy Measurement in Particle
        Physics}, 2nd ed.\ Oxford University Press, 2017.
  \item ATLAS Collaboration, Observation of a new boson at 125\,GeV,
        \textit{Phys.\ Lett.\ B} \textbf{716}, 1 (2012).
  \item FCC Collaboration, FCC-ee CDR Vol.\,2,
        \textit{Eur.\ Phys.\ J.\ ST} \textbf{228}, 261 (2019).
  \item M.-J.~Papa-Kango, Particle Shower Formation Simulations using
        Geant4, C++ and Python, simulation report, 2026.
  \item CERN Beamline for Schools.
        \url{https://beamlineforschools.cern}
\end{enumerate}

\vspace{12pt}
{\color{cernblue}\rule{\linewidth}{1pt}}
\vspace{4pt}

\noindent\small\itshape\color{gray}
\textbf{AI Use Declaration.}
AI was used for document structure only.

\end{document}
